% =============================================================================
%  Part IV: The Global T-Matrix  (Sections 19--22)
% =============================================================================

\part*{Part IV: The Global T-Matrix}

\section{From Local Site T-Matrix to Global Multiple Scattering}

The T-matrix $\bT_0$ derived in Parts~I and~II is a \emph{local site} (or
single-scatterer) quantity: it describes the self-consistent response of one
cubic cell in isolation, embedded in a homogeneous reference medium.  To
model a heterogeneous medium composed of many such cells, we must account
for the multiple scattering between sites---each cube scatters incident
waves, which propagate to neighbouring cubes, scatter again, and so on.

The standard framework for this is the Foldy--Lax multiple-scattering
theory \citep{Foldy1945,Lax1952}, recast in the T-matrix formalism by
\citet{Waterman1976} and applied to elastic media by Zeller \&
\citet{ZellerDederichs1973},
\citet{Gubernatis1977}, and \citet{Jakobsen2012}.

\subsection{The Lippmann--Schwinger Equation on a Lattice}

Partition the medium into $N$ non-overlapping cubes centred at positions
$\{\bx_n\}_{n=1}^N$, each with its own material contrasts
$\Delta c^{(n)}$, $\Delta\rho^{(n)}$.  The total scattered field at site~$m$
due to all other sites is:
\begin{equation}
  u_i^{\text{tot}}(\bx_m)
  = u_i^0(\bx_m)
  + \sum_{n\neq m}
    G_{ij}^{(0)}(\bx_m - \bx_n)\;\bT_0^{(n)}\cdot u_j^{\text{tot}}(\bx_n),
\label{eq:latticeLS}
\end{equation}
where $G_{ij}^{(0)}(\bx_m - \bx_n)$ is the background (reference medium)
Green's tensor evaluated at the vector separation between the centres of
cubes~$m$ and~$n$, and $\bT_0^{(n)}$ is the local site T-matrix for the
$n$-th cube (computed as in Part~II).

This is the elastic analogue of the Foldy--Lax
equations \citep{Foldy1945,Lax1952}.  In compact operator notation, writing
$\bu = [u(\bx_1),\ldots,u(\bx_N)]^T$ and assembling the inter-site
Green's tensors into a matrix~$\bG_0$:
\begin{equation}
  \bu = \bu^0 + \bG_0\,\bT_0\,\bu,
\label{eq:compactLS}
\end{equation}
where $\bT_0 = \mathrm{diag}(\bT_0^{(1)},\ldots,\bT_0^{(N)})$ is the
block-diagonal matrix of local site T-matrices.

\subsection{The Global T-Matrix}

Solving~\eqref{eq:compactLS} for~$\bu$ gives
$\bu = (\bI - \bG_0\,\bT_0)^{-1}\,\bu^0$.  The global T-matrix is
defined as the operator that maps the incident field to the total scattered
field:
\begin{equation}
  \bT = \bT_0\,(\bI - \bG_0\,\bT_0)^{-1}.
\label{eq:Tglobal}
\end{equation}
Equivalently, by pre-multiplying with $\bG_0$:
\begin{equation}
  \bG_0\,\bT
  = \bG_0\,\bT_0\,(\bI - \bG_0\,\bT_0)^{-1}
  = (\bI - \bG_0\,\bT_0)^{-1} - \bI.
\label{eq:G0T}
\end{equation}

\subsection{Physical Interpretation: The Neumann Series}

When $\|\bG_0\,\bT_0\| < 1$ (weak scattering), the inverse can be
expanded as a Neumann series:
\begin{equation}
  \bT = \bT_0 + \bT_0\,\bG_0\,\bT_0
  + \bT_0\,\bG_0\,\bT_0\,\bG_0\,\bT_0 + \cdots
\label{eq:Neumann}
\end{equation}

Each term has a transparent physical interpretation (Figure~1):

\begin{itemize}
\item \textbf{Zeroth order} $\bT_0$: Single scattering at each site (the
  Born--T-matrix approximation).

\item \textbf{First order} $\bT_0\,\bG_0\,\bT_0$: A wave scatters at
  site~$n$ ($\bT_0^{(n)}$), propagates to site~$m$ ($\bG_0$), and scatters
  again ($\bT_0^{(m)}$).

\item \textbf{$k$-th order:}  Chains of $(k+1)$ scattering events
  connected by $k$ propagation steps.
\end{itemize}
The full inverse $(\bI - \bG_0\,\bT_0)^{-1}$ sums this series to all
orders, capturing the complete multiple-scattering response.

\begin{quote}
\textit{[Schematic to be added.]}  A pictorial representation showing:
(a)~a regular lattice of cubes, each labelled with its local T-matrix
$\bT_0^{(n)}$;
(b)~arrows representing the inter-site Green's tensor
$G_{ij}^{(0)}(\bx_m-\bx_n)$ connecting cube centres;
(c)~the multiple-scattering chain: incident wave $\to$ scatter at site~$n$
$\to$ propagate $\to$ scatter at site~$m$ $\to$ propagate $\to\cdots$
\end{quote}

\noindent\textbf{Figure~1:}  Schematic of the global T-matrix formalism.
Left: the medium is partitioned into cubes, each with a local T-matrix
$\bT_0^{(n)}$.  Right: the inter-site Green's tensor $\bG_0$ propagates the
scattered field between cube centres, and the global T-matrix~\eqref{eq:Tglobal}
sums all multiple-scattering paths to all orders.


\section{The Inter-Site Green's Tensor $\bG_0$}

\subsection{Homogeneous Reference Medium}

In a homogeneous isotropic reference medium, the inter-site Green's tensor
is the free-space Green's tensor evaluated at the separation vector:
\begin{equation}
  G_{ij}^{(0)}(\bx_m - \bx_n)
  = f(|\bx_m-\bx_n|)\,\delta_{ij}
  + g(|\bx_m-\bx_n|)\,\hat{r}_{mn,i}\,\hat{r}_{mn,j},
\label{eq:G0homog}
\end{equation}
where $f(r)$ and $g(r)$ are the radial functions from Section~3, and
$\hat{r}_{mn,i} = (x_{m,i}-x_{n,i})/|\bx_m-\bx_n|$ is the unit vector
from cube~$n$ to cube~$m$.

For a regular cubic lattice with spacing $2a$ (so that adjacent cube faces
are touching), the inter-site vectors are
$\bx_m - \bx_n = 2a\,(p,q,r)$ for integers $(p,q,r)$ with
$(p,q,r)\neq(0,0,0)$.

\subsection{Layered Reference Medium: The Kennett Propagator}

For geophysical applications, the background medium is not homogeneous but
\emph{stratified}---consisting of horizontal layers with depth-dependent
properties.  In this case, the inter-site Green's tensor must be replaced
by the Green's tensor for the layered reference medium.

\citet{Kennett1983} showed that the elastic Green's tensor for a
plane-layered medium can be computed efficiently using the reflectivity
method: the wavefield is decomposed into plane waves (via horizontal
wavenumber integration), each of which propagates through the layer stack
via a system of reflection and transmission matrices.  The key recursive
relations are:

\begin{enumerate}
\item \textbf{Layer propagator:}  For a homogeneous layer of thickness~$h$
  with P- and S-wave vertical slownesses $q_\alpha$ and~$q_\beta$, the
  phase propagation across the layer is given by
  $\mathbf{E}(h) = \mathrm{diag}(e^{i\omega q_\alpha h},
  e^{i\omega q_\beta h})$.

\item \textbf{Interface coefficients:}  At each interface, Zoeppritz-type
  reflection and transmission coefficients $r_d, r_u, t_d, t_u$ (for both
  P--P, P--SV, SV--P, SV--SV conversions) relate the up- and down-going
  wave amplitudes.

\item \textbf{Kennett recursive algorithm:}  The generalised reflection
  matrix $R_D$ for a stack of layers is built recursively from the bottom
  up:
\begin{equation}
  R_D^{(j)} = r_d^{(j)}
  + t_u^{(j)}\,R_D^{(j+1)}\,
    \bigl(\bI - r_u^{(j)}\,R_D^{(j+1)}\bigr)^{-1}\,t_d^{(j)},
\label{eq:Kennett}
\end{equation}
where $r_d^{(j)}, t_d^{(j)}$ are the reflection/transmission matrices at
the $j$-th interface, and $R_D^{(j+1)}$ is the already-computed reflection
matrix for all layers below interface~$j$ \citep{Kennett1983,Kennett1979}.

\item \textbf{Green's tensor assembly:}  The frequency--wavenumber-domain
  Green's tensor is assembled from the reflection matrices and layer
  propagators.  A final inverse Hankel (or Fourier) transform over
  horizontal wavenumber yields
  $G_{ij}^{(0)}(\bx_m-\bx_n;\omega)$ in the space-frequency domain.
\end{enumerate}

\subsection{Significance for Heterogeneous Media Modelling}

The combination of the Kennett propagator for $\bG_0$ with the cubic local
T-matrix $\bT_0$ provides a powerful framework for modelling wave
propagation in realistic heterogeneous Earth models:

\begin{enumerate}
\item \textbf{Background model:}  A plane-layered velocity-density profile
  (e.g.\ from well logs or a 1D Earth model such as PREM or AK135) defines
  the reference medium.  The Green's tensor $\bG_0$ for this layered
  background is computed using the Kennett reflectivity method.

\item \textbf{3D heterogeneity:}  The departures from the layered background
  are partitioned into cubic cells, each with constant (but arbitrary)
  material contrasts $\Delta\lambda^{(n)}$, $\Delta\mu^{(n)}$,
  $\Delta\rho^{(n)}$.  The local T-matrix $\bT_0^{(n)}$ for each cube is
  computed analytically (Part~II)---it depends only on the cube's material
  contrasts and the local background properties.

\item \textbf{Multiple scattering:}  The global
  T-matrix~\eqref{eq:Tglobal} is assembled from the $\bT_0^{(n)}$ and
  $\bG_0$, and the total scattered field is obtained by applying $\bT$ to
  the incident field.
\end{enumerate}
This is essentially the approach of \citet{Jakobsen2012} and
\citet{JakobsenWu2016}, but with the important refinement
that the local T-matrix is computed self-consistently (not in the Born
approximation), allowing strong contrasts to be modelled accurately.


\section{Operator Structure and Symmetries}

\subsection{Block Structure of the Global System}

The global system~\eqref{eq:compactLS} has a natural block structure.  For
$N$ cubes, each described by 9 degrees of freedom (3 displacement + 6
strain components), the system matrices are of size $9N\times 9N$.  However,
the local T-matrix $\bT_0$ is block-diagonal (each block is $9\times 9$),
and the inter-site Green's tensor $\bG_0$ has translational structure (for a
regular lattice, $G_{mn}$ depends only on $\bx_m - \bx_n$).

For a regular cubic lattice, this translational invariance can be exploited:
the system~\eqref{eq:compactLS} becomes block-circulant and can be
diagonalised by a spatial Fourier transform over the lattice.  In Fourier
space, the global T-matrix at wavevector~$\mathbf{k}$ reduces to:
\begin{equation}
  \hat{\bT}(\mathbf{k})
  = \hat{\bT}_0\,\bigl(\bI - \hat{\bG}_0(\mathbf{k})\,\hat{\bT}_0\bigr)^{-1},
\label{eq:Tfourier}
\end{equation}
where $\hat{\bG}_0(\mathbf{k})
= \sum_{\bx_n\neq\mathbf{0}} \bG_0(\bx_n)\,e^{-i\mathbf{k}\cdot\bx_n}$
is the lattice Green's function (also called the structure factor of the
Green's tensor).

This Fourier decomposition reduces the $9N\times 9N$ system to $N$
independent $9\times 9$ systems, one for each wavevector in the first
Brillouin zone of the cubic lattice---a dramatic computational saving.

\subsection{Relationship to the Dyson Equation}

Equation~\eqref{eq:Tglobal} is the lattice analogue of the Dyson equation
from many-body theory.  Defining the self-energy (or mass operator) as
$\Sigma = \bT_0$, the effective Green's function of the heterogeneous
medium satisfies:
\begin{equation}
  \bG_{\text{eff}} = \bG_0 + \bG_0\,\Sigma\,\bG_{\text{eff}}
  = (\bI - \bG_0\,\Sigma)^{-1}\,\bG_0.
\label{eq:Dyson}
\end{equation}
The poles of $\bG_{\text{eff}}$ in the frequency--wavenumber domain give the
dispersion relations of the effective medium, including attenuation due to
scattering.  This connection to the Dyson equation was exploited by Zeller
\& \citet{ZellerDederichs1973} for elastic media and by
\citet{Jakobsen2012} for seismic modelling.


\section{Computational Methods for Inverting the Global System}

The central computational problem is to solve the linear system
\begin{equation}
  (\bI - \bG_0\,\bT_0)\,\bu = \bu^0,
\label{eq:globalsys}
\end{equation}
which is of dimension $9N\times 9N$ (three displacement and six independent
strain components per cube, $N$~cubes).  For realistic 3D models, $N$ may
range from $10^3$ (small exploration-scale problems) to $10^6$ or beyond
(crustal-scale), making the choice of solution method critical.

We survey five approaches, ranging from the simplest (direct inversion, the
Neumann series) to the most sophisticated (fast multipole methods).  Each
has a distinct domain of applicability.

\subsection{Direct Inversion}

The most straightforward approach is to form and factorise
$(\bI - \bG_0\,\bT_0)$ explicitly, using LU decomposition or a similar
dense solver.  This costs $O(N^3)$ in time and $O(N^2)$ in storage,
restricting it to small problems ($N\lesssim 10^3$).  However, it yields
the complete resolvent, from which the Green's function for all
source--receiver pairs can be extracted.  For benchmark problems and
validation of approximate methods, direct inversion remains indispensable.

\subsection{The Neumann (Born) Series}

The Neumann series~\eqref{eq:Neumann} is a fixed-point iteration:
\begin{equation}
  \bu^{(k+1)} = \bu^0 + \bG_0\,\bT_0\,\bu^{(k)}, \qquad
  \bu^{(0)} = \bu^0.
\label{eq:NeumannIter}
\end{equation}
It converges if and only if the spectral radius
$\rho(\bG_0\,\bT_0) < 1$.  In scattering theory this is the Born series,
and it is well known to diverge for strong material
contrasts \citep{Gubernatis1977}.  Even when it converges, the rate is
geometric with ratio~$\rho$, so many iterations may be needed for moderate
contrasts.  The Neumann series is therefore useful mainly as a conceptual
tool and as the starting point for more sophisticated iterative methods.

\subsection{Iterative Krylov Subspace Methods}

Krylov subspace methods---GMRES \citep{SaadSchultz1986},
BiCGSTAB \citep{vanderVorst1992}, and their variants---are the practical
workhorse for solving~\eqref{eq:globalsys}.  They require only the
matrix--vector product $\bv \mapsto \bG_0\,\bT_0\,\bv$, never the explicit
matrix, and achieve convergence in far fewer iterations than the Neumann
series because they minimise the residual over the full Krylov subspace
rather than iterating a single fixed-point map.

\textbf{GMRES} (Generalised Minimal Residual) \citep{SaadSchultz1986} is
the method of choice for non-symmetric systems.  At iteration~$k$ it
minimises $\|\mathbf{r}^{(k)}\|$ over the Krylov subspace
$\mathcal{K}_k = \mathrm{span}\{\mathbf{r}^{(0)},
\mathbf{A}\mathbf{r}^{(0)}, \ldots,
\mathbf{A}^{k-1}\mathbf{r}^{(0)}\}$, where
$\mathbf{A} = \bI - \bG_0\,\bT_0$.  Storage grows linearly with~$k$
(one vector per iteration), which can be controlled by restarting
(GMRES($m$)).

\textbf{BiCGSTAB} \citep{vanderVorst1992} is an alternative that avoids the
growing storage of GMRES at the cost of a less smooth convergence history.
Each iteration requires two matrix--vector products but only a fixed amount
of storage.

The matrix--vector product $\bG_0\,\bT_0\,\bv$ has a transparent physical
interpretation: (i)~\emph{scatter}: apply the block-diagonal $\bT_0$
to~$\bv$ (cost $O(N)$, since each $9\times 9$ block acts independently);
(ii)~\emph{propagate}: apply the dense inter-site Green's tensor $\bG_0$
(cost $O(N^2)$ naively, but reducible to $O(N\log N)$ with the methods of
Section~22.5).

\textbf{Preconditioning.}  The convergence rate of Krylov methods depends
on the spectral properties of $\bI - \bG_0\,\bT_0$.  Because the local
T-matrix $\bT_0$ already absorbs the strong single-site scattering
self-consistently, the remaining inter-site coupling $\bG_0\,\bT_0$ is
typically much weaker than the original material contrast, leading to
favourable convergence even without explicit preconditioning.  When
additional acceleration is needed, block-diagonal or block-Jacobi
preconditioners (using the diagonal blocks of $\bG_0$) are natural choices.

\textbf{Historical note.}  Krylov methods were applied to elastic-wave
scattering in layered media with superimposed spherical heterogeneities by
\citet{Nestor1996}, who solved the multiple-scattering series for the
global T-matrix using iterative methods with the Kennett reflectivity
method (Section~20.2) providing the inter-site Green's tensor $\bG_0$ and
spherical scatterer T-matrices for~$\bT_0$.  This combination demonstrated
the feasibility of modelling 3D heterogeneous media as perturbations to a
stratified background, but was limited by the isotropy of the spherical
T-matrix and the inability of spheres to tile space.  The cubic T-matrix
developed in Part~II removes both limitations, completing the bridge from
single-scatterer physics to rigorous homogenisation that was implicit in
the original formulation.

\subsection{The Renormalised Scattering Series}

\citet{JakobsenWu2016} proposed a renormalised scattering
series that overcomes the convergence limitations of the Born/Neumann
series for strong contrasts.  The idea is to split the total perturbation
into a ``strong'' part (absorbed into a modified reference medium or into
the local T-matrix) and a ``weak'' residual, then iterate on the residual
only.

Concretely, define a modified operator
$\tilde{\bG}_0 = \bG_0 - \bG_0^{\text{diag}}$, where
$\bG_0^{\text{diag}}$ contains only the self-interaction terms (which are
already accounted for in $\bT_0$).  The renormalised series iterates:
\begin{equation}
  \bu^{(k+1)} = \bu^0 + \tilde{\bG}_0\,\bT_0\,\bu^{(k)}.
\label{eq:renorm}
\end{equation}
Because $\tilde{\bG}_0$ excludes the (already-resummed) on-site terms, the
spectral radius $\rho(\tilde{\bG}_0\,\bT_0)$ is typically much smaller
than $\rho(\bG_0\,\bT_0)$, extending the domain of convergence to stronger
contrasts.  When combined with the self-consistent (non-Born) local
T-matrix, this approach effectively sums the dominant scattering
contributions exactly and treats only the residual inter-site coupling
perturbatively.

The renormalised series can also be viewed as a preconditioned Neumann
iteration, and it can be further accelerated by embedding it within a
Krylov solver.

\subsection{Fast Multipole and Hierarchical Matrix Methods}

The dominant cost of each Krylov iteration is the dense matrix--vector
product $\bG_0\,\bv$, which costs $O(N^2)$ when computed directly.  For
large~$N$, this can be reduced to $O(N\log N)$ or even $O(N)$ by
exploiting the spatial structure of the Green's tensor.

\textbf{Fast Multipole Method (FMM):}  The elastic Green's tensor
$G_{ij}^{(0)}(\bx)$ decays as $1/|\bx|$ and is smooth (low-rank) when
source and receiver are well separated.  The
FMM \citep{Rokhlin1985,GreengardRokhlin1987} exploits this by grouping
cubes into a hierarchical tree structure, computing far-field interactions
via multipole expansions, and only evaluating near-field interactions
directly.  For elastodynamics, the FMM was developed by
\citet{Fujiwara2000} and \citet{Chaillat2008},
achieving $O(N\log N)$ or $O(N)$ complexity per matrix--vector product.

\textbf{Hierarchical matrices ($\Hmat$-matrices):}  An alternative to the
FMM, $\Hmat$-matrices \citep{Hackbusch1999,Boerm2003} represent the
off-diagonal blocks of $\bG_0$ in low-rank factored form, again achieving
$O(N\log N)$ storage and matrix--vector products.  The advantage over FMM
is that $\Hmat$-matrices can also be inverted or LU-factored in
$O(N\log^2 N)$ time, providing a direct (non-iterative) solver.

\textbf{Hybrid FFT for layered $\bG_0$:}  When the reference medium is
layered (Section~20.2) and the cubes are arranged on a regular horizontal
grid (but at variable depths), the horizontal convolution structure of
$\bG_0$ can be exploited via 2D FFT.  Only the vertical coupling between
different depth levels requires explicit summation.  This hybrid approach
naturally suits the ``layered background + 3D perturbations'' geometry of
geophysical applications.


\subsection{Summary and Recommendations}

Table~1 summarises the computational methods and their domains of
applicability.

For the geophysical application envisioned here---3D heterogeneities
superimposed on a layered background via the Kennett propagator---the
recommended approach is a Krylov solver (GMRES or BiCGSTAB) with the
matrix--vector product $\bG_0\,\bT_0\,\bv$ accelerated by the hybrid
FFT/direct scheme described in Section~22.5.  The self-consistent local
T-matrix $\bT_0$ from Part~II serves as a natural preconditioner by
absorbing the strong single-site scattering, leaving only the weaker
inter-site coupling to be resolved iteratively.

\begin{table}[ht]
\centering
\begin{tabular}{llll}
\toprule
\textbf{Method} & \textbf{Cost} & \textbf{Requires} & \textbf{Best for} \\
\midrule
Direct inversion      & $O(N^3)$                       & Dense matrix              & Small $N$, benchmarks \\
Neumann series        & $O(kN^2)$                      & $\rho < 1$                & Weak scattering only \\
GMRES/BiCGSTAB        & $O(kN^2)$ or $O(kN\!\log\!N)$  & Mat-vec product           & General heterogeneity \\
Renormalised series   & $O(kN^2)$ or $O(kN\!\log\!N)$  & Self-consist.\ $T_0$      & Strong contrasts \\
FMM / $\Hmat$-matrices & $O(N\!\log\!N)$ per iter. & Spatial decay of $G_0$ & Large $N$, 3D \\
\bottomrule
\end{tabular}
\caption{Comparison of computational methods for solving the global T-matrix system~\eqref{eq:globalsys}.
Here $k$ denotes the number of iterations required for convergence.}
\end{table}

\bigskip
\noindent
\textit{Generated from \texttt{tmatrix\_analytic.py} and \texttt{tmatrix\_cube.py}
using SymPy symbolic computation.  All sphere integrals evaluated in exact closed
form; cube integrals evaluated semi-analytically via Taylor expansion to
$O(\omega a)^8$.}
